% ==============================================
% CFD Tutorial Template: OpenFOAM Terminal Basics
% Class: scrartcl (KOMA-Script) - optimized for tutorials
% Content: pitzDaily with simpleFoam
% ==============================================

\documentclass[
    parskip=half,      % Space between paragraphs (better for steps)
    DIV=12,            % Optimal margins for readability
    headings=normal,   % Professional heading sizes
    fontsize=11pt,     % Slightly larger for screen reading
    english            % Language setting
]{scrartcl}

\input{../preamble.tex}

% ========== DOCUMENT START ==========
\begin{document}

% ========== TITLE PAGE ==========
\title{Tutorial \#2: Water-air Ejector}
\subtitle{Multiphase Transient Flow}
\subject{Workshop on Computational Fluid Dynamics (CFD) \\ with OpenFOAM \\ IMP-PAN Gdansk}
\author{Robert Castilla \\ CATMech - Universitat Politècnica de Catalunya \\ Department of Fluid Mechanics}
\date{2025-26 \\ Version 1.0}



\maketitle

% ========== ABSTRACT ==========
\begin{abstract}

\end{abstract}

\tableofcontents
\newpage

\fbox{
	\parbox{15cm}{

\section*{Learning Objectives}
By the end of this tutorial the student will be able to
	
	\begin{enumerate}
		\item \textbf{Navigate} OpenFOAM case directory structure
		
		\item \textbf{Understand} basic OpenFOAM file formats
		
		\item \textbf{Run} a simple transient simulation
		
		\item \textbf{Analyze} basic results in terminal and \terminalcmd{paraview}
	\end{enumerate}
}
}
% ========== SECTION 1: INTRODUCTION ==========
\section{Introduction and Theoretical Background}


\subsection{Case Overview}

% ========== SECTION 2: TERMINAL SETUP ==========
\section{Case Setup}

\subsection{Python virtual environment}




% ========== SECTION 4: RUNNING THE SIMULATION ==========
\section{Running the Simulation}


% ========== SECTION 5: POST-PROCESSING ==========
\section{Post-Processing in Terminal}

\section{Visualization in paraview}

% ========== SECTION 7: EXERCISES ==========
\section{Student Exercises and Extensions}


% ========== DISCLAIMER ============

\section*{Disclaimer}

This tutorial has been developed by \textbf{Robert Castilla}, 
of the Departament of Fluid
Mechanics in the Universitat Politècnica de Catalunya, for the doctoral 
workshop at \textbf{IMP⁻PAN Gda\'{n}sk} (July 2026). The author utilized \textbf{DeepSeek} for the initial
definition of the document format. Further technical refinement and the development
of advanced post-processing exercises were supported by \textbf{NotebookLM}. All the technical
information for the \textbf{OpenFOAM v2512} environmnet and all physical interpretations were
performed and approved by the author. 

% ========== END DOCUMENT ==========
\end{document}